\documentclass[11pt,a4paper]{article}
\usepackage[utf8]{inputenc}
\usepackage[spanish]{babel}
\usepackage{geometry}
\geometry{top=2.5cm, bottom=2.5cm, left=2.5cm, right=2.5cm}
\usepackage{amsmath,amssymb}
\usepackage{booktabs}
\usepackage{tabularx}
\usepackage{xcolor}
\usepackage{tcolorbox}
\usepackage{hyperref}
\usepackage{enumitem}

\definecolor{primary}{RGB}{15, 32, 67}
\definecolor{secondary}{RGB}{0, 114, 206}
\definecolor{lightgray}{RGB}{245, 247, 250}

\hypersetup{
    colorlinks=true,
    linkcolor=secondary,
    urlcolor=secondary,
    pdftitle={Tarea 1 - Validacion del Problema con Expertos}
}

\title{
    \vspace{-1cm}
    \textbf{\huge Tarea 1: Taller de Ideación y Validación con Expertos}\\[0.3cm]
    \large\textbf{Proyecto StatCredit AI}\\[0.1cm]
    \small Universidad Externado de Colombia -- Emprendimiento (Ciencia de Datos)\\[0.1cm]
    \small\textbf{Grupo 5: Santiago Sandoval | Sebastián Ramos | Julián Duarte | Tomás Rincón | Julián Jiménez}
}
\author{}
\date{Semestre 2026-II}

\begin{document}

\maketitle

\begin{tcolorbox}[colback=primary!5,colframe=primary,title=\textbf{¿De qué trata este entregable?}]
En esta primera tarea validamos en campo la hipótesis del problema de \textbf{StatCredit AI} mediante entrevistas cualitativas en profundidad con expertos del sector financiero en Colombia (un ejecutivo de riesgo crediticio de consumo y una experta en regulación actuarial SARC). El objetivo es corroborar si el punto ciego de los independientes y la necesidad de explicabilidad en Inteligencia Artificial representan dolores reales y vigentes en la banca.
\end{tcolorbox}

\section{El Problema y Reencuadre Metodológico}

Muchos trabajadores independientes solventes en Colombia (comerciantes urbanos, profesionales autónomos y micronegociantes) poseen una intensa actividad financiera diaria en billeteras digitales (Nequi, Daviplata) y facturación electrónica. Sin embargo, al solicitar crédito formal, son declinados por carecer de historial crediticio bancario o certificaciones tradicionales de nómina.

\textbf{Nuestro reencuadre del problema (Bianchi \& Verganti, 2021):}
El problema no es que los trabajadores independientes carezcan de capacidad de repago. El verdadero problema es que los procesos tradicionales de evaluación no cuentan con la infraestructura analítica para ingestar e interpretar esa información transaccional observable. En StatCredit AI no buscamos cambiar a las personas, sino cambiar la forma en que se mide quién puede pagar un crédito.

\section{Caracterización de los Expertos Entrevistados}

Para validar estas hipótesis, realizamos entrevistas cualitativas estructuradas a dos perfiles estratégicos:

\begin{enumerate}[leftmargin=*]
    \item \textbf{Diego Ramos (Experto en Riesgo Crediticio):} Ejecutivo de crédito de consumo con años de experiencia directa en la evaluación operacional y comités de aprobación crediticia.
    \item \textbf{Paola Cárdenas (Experta en Regulación SARC):} Analista de crédito microempresarial y especialista en cumplimiento normativo de la Superintendencia Financiera de Colombia (SARC).
\end{enumerate}

\section{Transcripción y Resultados de las Entrevistas de Campo}

\subsection{Entrevista 1: Diego Ramos (Ejecutivo de Crédito de Consumo)}

\begin{tcolorbox}[colback=lightgray,colframe=primary,title=\textbf{Pregunta 1: Rechazo de independientes por historial en buró}]
\textit{¿Ustedes rechazan mucha gente independiente por lo del tema de DataCrédito y esas cosas?}
\tcblower
\textbf{Respuesta de Diego Ramos:}  
«Sí, muchísimo, más de lo que se cree normalmente. Yo llevaba años en eso y apenas el sistema ve que el cliente no tiene un puntaje bueno en el buró, ni siquiera llega a que un analista lo mire con calma. Y muchas veces ese señor sí tiene con qué pagar por lo que puede decirnos más personalmente, solo que nunca tuvo una tarjeta de crédito o nunca sacó un préstamo formal, entonces para el sistema es como si no existiera.»
\end{tcolorbox}

\begin{tcolorbox}[colback=lightgray,colframe=primary,title=\textbf{Pregunta 2: Capacidad de procesamiento manual de extractos digitales}]
\textit{¿Y si la persona le muestra los movimientos de Nequi o las facturas, ustedes tienen cómo revisar eso rápido?}
\tcblower
\textbf{Respuesta de Diego Ramos:}  
«No, y ahí es donde hay un problema ya que eso tocaría revisarlo a mano, uno abriendo el PDF del extracto, contando transacción por transacción, sacando el promedio en Excel y así uno alcanza a sacar máximo 15 o 20 casos al día por persona. Pero no es que no queramos usar esa información, es que no damos abasto para procesarla cuando hay tantas personas que manejan el dinero de esa manera. Si a mí me hubieran dado algo automático que me diera bien la información como lo que me cuentas que quieren hacer, eso me cambiaría la manera de cómo se rechazaría a los clientes.»
\end{tcolorbox}

\begin{tcolorbox}[colback=lightgray,colframe=primary,title=\textbf{Pregunta 3: Comportamiento de pago de independientes vs. asalariados}]
\textit{¿Un independiente que mueve plata todos los días es más moroso que uno con sueldo fijo?}
\tcblower
\textbf{Respuesta de Diego Ramos:}  
«No, y digamos que eso lo sabemos todos los que trabajan en esto, aunque no se mencione mucho. Yo he visto casos de gente que tiene su negocio como yo y le entra plata todos los días, y hay más organizados que muchos asalariados que conozco. Lo que pasa es que a uno en el banco lo criaron con el cuento de que 'nómina fija es igual a seguridad', y es como que ese chip es difícil de cambiar para que la gente de un día para otro revise otras perspectivas.»
\end{tcolorbox}

\begin{tcolorbox}[colback=lightgray,colframe=primary,title=\textbf{Pregunta 4: Miedo a la aprobación sin desprendible de nómina}]
\textit{¿Y qué tanto miedo le tienen ustedes a aprobar un crédito sin el desprendible de nómina de toda la vida?}
\tcblower
\textbf{Respuesta de Diego Ramos:}  
«El miedo no es nuestro sino del banco, porque muchas personas nos cuentan sus historias y de manera hablada catalogan para poder entrar a revisión pero sin las herramientas deja de ser posible. Si yo apruebo algo raro y después algo sale mal, quien tiene que ir a dar la cara en el comité de auditoría soy yo. Entonces pues uno prefiere por si las moscas, decir que no de una vez, así uno sepa que se le está negando el préstamo a alguien que seguramente sí paga.»
\end{tcolorbox}

\subsection{Entrevista 2: Paola Cárdenas (Analista de Crédito Microempresarial \& Experta SARC)}

\begin{tcolorbox}[colback=lightgray,colframe=secondary,title=\textbf{Pregunta 5: Requerimientos de la Superintendencia Financiera (SARC)}]
\textit{¿Qué les exige la Superfinanciera cuando un banco quiere usar algoritmos para dar créditos?}
\tcblower
\textbf{Respuesta de Paola Cárdenas:}  
«Lo que a uno le exigen de fondo con el tema del SARC es que la entidad pueda explicar en cualquier momento por qué se tomó tal decisión de crédito. Entonces a uno no le sirve decir 'es que el modelo dijo que no' y ya. Da igual si es una fórmula de estadística o un algoritmo de inteligencia artificial, la responsabilidad de explicarlo siempre es de la entidad, nunca del programa ni de quien se lo vendió.»
\end{tcolorbox}

\begin{tcolorbox}[colback=lightgray,colframe=secondary,title=\textbf{Pregunta 6: Tranquilidad regulatoria con reportes de explicabilidad (XAI/SHAP)}]
\textit{¿Y si nuestra herramienta les entrega un reporte fácil de entender explicando por qué la IA decidió cada cosa, eso les da tranquilidad para una auditoría?}
\tcblower
\textbf{Respuesta de Paola Cárdenas:}  
«Sí, digamos que si tienen un pleno soporte que el banco acepte se podría hacer ya que nos autorizarían a poder utilizarlo como un complemento, pero sin ese reportecito que explique qué pesó a favor y qué en contra, ningún oficial de cumplimiento se va a poner a defender esa decisión frente a un auditor, así el modelo sea buenísimo o ustedes tengan mucha información. Así que con esa explicación clara, la cosa deja de sonar a caja negra sospechosa y uno sí la puede defender más tranquilo, o en caso tal no toda la responsabilidad cae sobre el que acepta el crédito.»
\end{tcolorbox}

\section{Síntesis de Hallazgos y Validación de Hipótesis}

Las entrevistas de campo aportan evidencias cualitativas directas que respaldan de forma contundente la propuesta de valor de **StatCredit AI**:

\begin{table}[h!]
\centering
\small
\begin{tabularx}{\textwidth}{p{3.5cm} X p{4.5cm}}
\toprule
\textbf{Hipótesis Evaluada} & \textbf{Evidencia Cualitativa Obtenida en Campo} & \textbf{Decisión de Diseño en StatCredit AI} \\
\midrule
\textbf{1. Fricción por Ingesta Manual} & Diego Ramos confirma que los analistas solo pueden procesar 15-20 extractos al día manualmente en Excel. & Desarrollo de la API REST para ingesta y cálculo automatizado en $<2$ segundos. \\
\midrule
\textbf{2. Falsos Negativos y Sesgo Cultural} & Diego evidencia que el sesgo de "nómina fija igual a seguridad" causa rechazo masivo de independientes solventes. & Implementación de modelos de Inferencia Causal para evaluar estabilidad en flujos estacionales. \\
\midrule
\textbf{3. Miedo al Comité y Explicabilidad SARC} & Paola Cárdenas señala que ningún oficial defenderá un modelo de "caja negra" ante la Superfinanciera (Circular 026). & Inclusión nativa de la matriz de explicabilidad SHAP en la respuesta JSON y el Dashboard. \\
\bottomrule
\end{tabularx}
\caption{Matriz de validación de hipótesis a partir de los testimonios de los expertos.}
\end{table}

\section{Conclusión}

La validación con los expertos Diego Ramos y Paola Cárdenas confirma que el dolor en el otorgamiento de crédito a independientes no se debe a desinterés institucional, sino a la falta de herramientas tecnológicas automatizadas de ingesta transaccional y reportes explicables adaptados a las exigencias del SARC. **StatCredit AI resuelve exactamente esta fricción operativa y regulatoria.**

\end{document}
